% Verified: True
% Notes to assembler: Preamble deps needed: \\usepackage{algorithm}, \\usepackage{algpseudocode}, \\usepackage{booktabs}, \\usepackage{amsthm} with a `definition` theorem style (Definition 4 = urgency, Definition 5 = skip filter — numbering is intentional per the prompt), \\usepackage{siunitx} (not used here but consiste

\subsection{Multi-Signal Urgency Triage}
\label{sec:urgency}

Once the store is populated and enriched with a \WorkGraph{} (\S\ref{sec:workgraph}),
the remaining problem is inbox-scale: how does the system decide what,
among $N$ heterogeneous records, actually deserves the user's attention
this hour?  \coviber{}'s answer is a rule-based additive scorer with a
prefilter, implemented in \path{coviber/urgency.py}. The design goal is
explicit: the scoring function must be inspectable, deterministic, and
reproducible on-device, with no reliance on a hosted classifier and no
labeled training corpus.

\subsubsection{Formal definition}

Let $\mathcal{S} = \{s_1, \ldots, s_{10}\}$ be the fixed set of signal
predicates enumerated in Table~\ref{tab:signals}. Each signal is a
Boolean predicate $s_i : \mathcal{R} \to \{0,1\}$ over a record $r$
(cf.\ Definition~\ref{def:record}). Let $w : \mathcal{S} \to
\mathbb{Z}_{\ge 0}$ be a configured weight map. The three age-tier
signals $\{s_{\text{age\_7d}}, s_{\text{age\_48h}}, s_{\text{age\_24h}}\}$
are mutually exclusive by construction (they are the branches of a
single \texttt{if/elif/elif} on record age), so at most one fires per
record.

\begin{definition}[Urgency score]
\label{def:urgency}
The urgency of a record $r$ is
\[
U(r) \;=\; \sum_{s \in S(r)} w(s), \qquad
S(r) \;=\; \{\, s \in \mathcal{S} \mid s(r) = 1 \,\},
\]
with a signal contributing zero (and being treated as disabled) iff
$w(s) \le 0$.
\end{definition}

At the default weights shipped with the package
(Table~\ref{tab:signals}), the seven non-age signals sum to $11$ and the
maximal age tier contributes $3$, so $U(r) \in [0, 14]$.  This
$[0,14]$ contract is pinned by a regression test
(\texttt{test\_urgency\_contract}); no clamp is applied at runtime.
Under a custom weight map the ceiling shifts in proportion to the
caller's choices, which is deliberate: silent clamping would mask a
misconfiguration rather than surface it. Configs that supply only a
partial weight dictionary are merged key-by-key against the defaults,
unknown keys are surfaced as \texttt{RuntimeWarning}, and non-integer
values raise a \texttt{ValueError} that names the offending key. Weight
$0$ is a first-class way to disable a signal.

\begin{definition}[Skip filter]
\label{def:skip}
A record $r$ is \emph{skipped} (excluded from the queue before scoring)
iff
\[
\mathrm{skip}(r) \;=\; \mathrm{skip\_sender}(r) \;\lor\; \mathrm{skip\_subject}(r) \;\lor\; \mathrm{fyi}(r),
\]
where $\mathrm{skip\_sender}$ matches any element of
$\mathtt{cfg.skip\_senders}$ against \texttt{r.from\_name} on token
boundaries, $\mathrm{skip\_subject}$ matches any element of
$\mathtt{cfg.skip\_subjects}$ as a case-insensitive substring of
\texttt{r.subject}, and $\mathrm{fyi}$ matches any element of the FYI
lexicon (defaults: \texttt{fyi}, \texttt{no action needed},
\texttt{no reply needed}, \texttt{just so you know}) against the
concatenation of subject and body, again on token boundaries.
\end{definition}

The token-boundary discipline for both $\mathrm{skip\_sender}$ and
$\mathrm{fyi}$ is load-bearing. Prior versions of \coviber{} matched
these lexicons as raw substrings, which caused two classes of false
positive: (i) colleagues whose surname literally contained a skipped
token (e.g. \texttt{"Abbott"} matching \texttt{"bot"}), and (ii)
common English words containing the trigram \texttt{"fyi"} inside them
(\texttt{"justifying"}, \texttt{"notifying"}, \texttt{"typify"}). Both
would silently drop legitimate obligations from the queue. The v0.4/L5
audit tightened both matches to word-boundary regexes of the form
\verb|(?<![a-z0-9])<lex>(?![a-z0-9])|; the fix is small, but it is the
canonical example of the transparency guarantee that motivates
rule-based scoring in the first place.

The complete triage procedure is given in Algorithm~\ref{alg:triage}.
Three additional filters run in the loop but are not part of
Definition~\ref{def:skip}: self-authored records (where
\texttt{r.from\_name} equals \texttt{cfg.you}) are dropped as
non-obligations, the skip filter runs first, and any record that scores
$U(r) \le 0$ after scoring is dropped as noise. The surviving records
are sorted by $(-U(r), r.\mathrm{ts})$ so that ties break in favour of
older records.

\begin{algorithm}
\caption{Triage: filter, score, sort.}
\label{alg:triage}
\begin{algorithmic}[1]
\Require records $R$, urgency config $\mathrm{cfg}$
\Ensure ranked queue $Q$ of $\{\mathrm{record}, U, S(r)\}$
\State $Q \gets [\,]$
\ForAll{$r \in R$}
  \If{$r.\mathrm{from\_name}.\mathrm{lower}() = \mathrm{cfg.you}.\mathrm{lower}()$} \State \textbf{continue} \EndIf
  \If{$\mathrm{skip}(r, \mathrm{cfg}) \neq \varnothing$} \State \textbf{continue} \EndIf
  \State $(u, S) \gets \mathrm{score}(r, \mathrm{cfg})$
  \If{$u \le 0$} \State \textbf{continue} \EndIf
  \State $Q.\mathrm{append}(\{\mathrm{record}\!:\!r,\, U\!:\!u,\, S\!:\!S\})$
\EndFor
\State \Return $Q$ sorted by $(-U, r.\mathrm{ts})$
\end{algorithmic}
\end{algorithm}

Complexity is $O(|R| \cdot k)$ per call, with $k$ the maximum lexicon
size (action words, priority senders, collaborators, skip lexicons).
The predicates are each $O(1)$ or an $O(k)$ substring/regex scan over a
bounded token blob; no cross-record work is done, so triage scales
linearly with the store and admits trivial parallelisation.

\subsubsection{The signal set}

\begin{table*}[t]
\centering
\small
\caption{The ten urgency signal predicates. Default weights sum to
$11$ (non-age) plus at most $3$ (age tier), giving the
$U(r) \in [0, 14]$ contract at defaults. Weights are per-key
overridable; unknown keys warn, non-integers raise.}
\label{tab:signals}
\begin{tabular}{lcp{9.2cm}}
\toprule
\textbf{Signal} & \textbf{Default $w$} & \textbf{Predicate / rationale} \\
\midrule
\texttt{mention}         & 3 & Regex \verb|@<you>\b| matches in \texttt{r.text} (case-insensitive). Empty \texttt{cfg.you} is guarded to avoid collapsing the regex to \verb|@\b|, which would fire on every raw \texttt{@}. Direct addressing is the single strongest human signal of obligation. \\
\texttt{priority\_sender} & 2 & \texttt{r.from\_name.lower()} $\in$ \texttt{cfg.priority\_senders}. Configured per-user (managers, VIPs). Off by default (empty set); nothing about any specific person or organization is baked in. \\
\texttt{action\_word}    & 2 & Any token from the action lexicon (defaults: \texttt{please}, \texttt{can you}, \texttt{need}, \texttt{review}, \texttt{approve}, \texttt{decision}, \texttt{deadline}, \texttt{asap}, \texttt{urgent}, \texttt{waiting on}, \texttt{action required}, \dots) appears in \texttt{subject $\|$ text}. Language-level cue that a reply is expected. \\
\texttt{question}        & 1 & \texttt{'?'} in \texttt{r.text}. Weak cue but cheap and orthogonal to action words. \\
\texttt{unread}          & 1 & \texttt{r.unread} is true. Loader-supplied; degrades gracefully to $0$ for sources with no read state. \\
\texttt{no\_reply}       & 1 & \texttt{r.thread\_id} is set and \texttt{r.replied} is false --- the user has not yet responded in a live thread. \\
\texttt{collaborator}    & 1 & \texttt{r.from\_name.lower()} $\in$ \texttt{cfg.collaborators}. Nudges known peers up without over-boosting them. \\
\texttt{age\_7d}         & 3 & Record age $> 168$\,h. Mutually exclusive with the finer age tiers. Escalates neglected obligations. \\
\texttt{age\_48h}        & 2 & Record age $\in (48, 168]$\,h. \\
\texttt{age\_24h}        & 1 & Record age $\in (24, 48]$\,h. \\
\bottomrule
\end{tabular}
\end{table*}

Two properties of the signal set deserve emphasis. First, the signals
are \emph{orthogonal by design}: \texttt{mention} is content-derived,
\texttt{priority\_sender} is identity-derived, \texttt{unread} is
platform-state-derived, and the age tiers are time-derived. A single
record is expected to fire two to four signals; the arithmetic sum is
then a joint score across independent evidence axes. Second, all
identity-bearing lexicons --- priority senders, collaborators, action
words, skip senders/subjects --- default either to an empty set or to a
generic vocabulary (mail-bot names, ``weekly digest'' subjects). No
individual person or organization is hard-coded. This is what makes the
scorer portable across users without a preprocessing step and safe to
ship as an open-source default.

\subsubsection{Design contrast: rule-based vs.\ learned}

The prevailing academic and industrial approach to inbox
prioritisation is a supervised classifier ---
\cite{TODO-horvitz-priorities} in the classical setting, and more
recently BERT-style text encoders or LLM zero-shot rankers fronted by a
managed API. Both classes buy accuracy at three costs \coviber{}
declines to pay:

\begin{enumerate}
\item \textbf{Labeled data.} A classifier needs per-user labels
  (\emph{this was urgent}, \emph{this was noise}) to be well-calibrated
  for a specific mailbox. Rule-based scoring calibrates via
  configuration, which the user writes once and can inspect.
\item \textbf{Cloud egress.} Zero-shot LLM ranking exports the entire
  inbox contents to a hosted model. \coviber{}'s target deployment is
  local-first over sensitive corpora (\S\ref{sec:threat-model}); this
  is a non-starter.
\item \textbf{Opacity.} A neural score is a scalar with no
  human-legible provenance. Definition~\ref{def:urgency} returns
  $(U(r), S(r))$: every triage decision ships with the list of signals
  that fired and their individual weights, which is what makes the FYI
  regression above visible in the first place.
\end{enumerate}

The tradeoff is real: rule-based scoring cannot learn user-specific
idioms without configuration edits, and it will underrank a
subtly-urgent record that fires zero signals. The system is honest
about this: the persona layer (\S\ref{sec:persona}) and the MCP recall
surface (\S\ref{sec:mcp}) are how a downstream LLM client re-ranks or
overrides triage output on demand. Urgency is the fast, transparent
first pass; the LLM is the slow, contextual second pass.

